\section{Conclusions}
\label{sec:main:conclusion}
Our study analyzes the training dynamics of two-layer neural networks for learning multi-index target functions, distinctively focusing on {\it multi-pass} gradient descent which involves reusing batches multiple times. We find that this enables gradient descent to exceed the constraints imposed by information and leap exponents. 

Gradient descent is found to achieve a positive correlation with the target function across a broader class than previously anticipated, with only two data batch repetitions. Our analysis further demonstrates that the limitations associated with information and leap exponents, staircase learning, and CSQ lower bounds are restricted to online/single pass SGD and do not describe the class of functions inherently easy or hard to learn by gradient-based methods for neural networks. 

Our conclusions follow from rigorous mathematical proofs derived from Dynamical Mean Field Theory, through which we also offer an analytical description of the dynamic processes of low-dimensional weight projections—a noteworthy insight. Additionally, we provide a closed-form depiction of these dynamical processes and illustrate our theoretical findings with numerical experiments.
